\origpage{51}{}
\noindent \textit{$T \setminus \{f_1, f_2\}$ where $f_i$ has ends in $T_0$ and $T_i$ ($i = 1, 2$). Suppose that there is a unique set of $k$ mutually vertex-disjoint paths of $W(T_0)$ between $W(f_1)$ and $W(f_2)$ and every vertex of $W(T_0)$ belongs to one of these paths. Then $V(T_0) \cap N(T) = \emptyset$.}

\textit{Subproof.} Let $A = W(T_1), B = W(T_0) \cup W(T_2), A' = W(T_0) \cup W(T_1), B' = W(T_2)$. Then $B \cap A' = W(T_0)$. Choose a path-decomposition $(P, W')$ as in (12.1) (with $W$ replaced by $W'$)\fix{, as built in the proof of (12.1); by the same induction, for each $t \in V(P)$ that is not an end of $P$, every vertex of $W'(t)$ belongs to $W'(e)$ for some edge $e$ of $P$ incident with $t$, as (12.2) needs}.\fixnote{fix:12.3-claim1}{This is the hypothesis added to (12.2). In the proof of (12.1), each of Claims (1) and (2) adds a new end $s$ next to $s_1$, and gives $s_1$ a graph with $V(W(s_1)) = V(A \cap B) \cup V(A_1 \cap B_1)$, which is the union of the sets $W(e)$ of the two edges of $P$ at $s_1$. The other inner vertices keep their graphs.} We may assume that $P$ has ends $s_1, s_2$, where $s_i$ is the end of $f_i$ in $T_i$ ($i = 1, 2$), and otherwise $P$ is disjoint from $T$. Let $(T^*, W^*)$ be as in (12.2). Then $N(T^*) \subseteq N(T)$, and so equality holds, from the choice of $(T, W)$; but $V(T_0) \cap N(T^*) = \emptyset$, and so $V(T_0) \cap N(T) = \emptyset$. This proves (1).

\hspace*{12pt}(2) \textit{$|V(P) \cap N(T)| \le n$ for every path $P$ of $T$.}

\textit{Subproof.} Suppose not; then there is a path $P$ of $T$ with both ends in $N(T)$ and with $|V(P) \cap N(T)| = n + 1$. Let $V(P) \cap N(T) = \{t_0, t_1, \dots, t_n\}$, in order on $P$. For $1 \le i \le n$, let $f_i \in E(P)$, chosen so that $f_i$ is between $t_{i-1}$ and $t_i$, and of all such edges $|W(f_i)|$ is minimum; let $S_i, T_i$ be the two components of $T \setminus \delm{e_i}\fix{f_i}$ where $t_0 \in V(S_i)$, and let $A_i = W(S_i), B_i = W(T_i)$. Then for $1 \le i < j \le n$, $A_i \subseteq A_j$ and $B_j \subseteq B_i$. Suppose that $1 \le i < i' \le n$, and $(A_i, B_i)$ and $(A_{i'}, B_{i'})$ have the same order $k$ say, and $(A_j, B_j)$ has order $> k$ for $i < j < i'$. In other words, $|W(f_i)| = |W(f_{i'})| = k$, and $|W(f_j)| > k$ for $i < j < i'$. From the definition of $f_j$ ($1 \le j \le n$), it follows that $|W(e)| \ge k$ for all edges $e$ of the path of $T$ between $f_i$ and $f_{i'}$. Since $(W, T)$ is linked, there are $k$ mutually vertex-disjoint paths of $G$ between $W(f_i)$ and $W(f_{i'})$, and so there is a linkage $M_{ii'}$ in $B_i \cap A_{i'}$ with $k$ components, each with one end in $V(A_i \cap B_i)$ and the other in $V(A_{i'} \cap B_{i'})$. By (1), since there is a member of $N(T)$ between $\delm{e_i}\fix{f_i}$ and $\delm{e_{i'}}\fix{f_{i'}}$, it follows that either $V(M_{ii'}) \ne V(B_i \cap A_{i'})$, or there is more than one choice for $M_{ii'}$. In either case, we may choose $M_{ii'}$ so that $M_{ii'} \ne L \cap B_i \cap A_{i'}$, since $V(L \cap B_i \cap A_{i'}) = V(B_i \cap A_{i'})$. But this contradicts (2.6), since each $(A_j, B_j)$ has order $\le \delm{w + 1 = k}\fix{w' + 1}$.\fixnote{fix:12.3-claim2}{The edges are called $f_i$, so ``$T \setminus e_i$'' and ``between $e_i$ and $e_{i'}$'' mean $T \setminus f_i$ and between $f_i$ and $f_{i'}$. The width of $(T, W)$ is at most $w'$, so the orders are at most $w' + 1$, which is the $k$ for which $n$ was chosen in (2.6). The $k$ of the claim is another number.} Hence there is no such path $P$, and so (2) holds.

Since $T$ has maximum degree $\le 3$, it follows from (2) that $|N(T)| \le 2^n$. \fix{To see this, for adjacent $u, v \in V(T)$ let the \textit{branch} at $u$ through $v$ be the set of vertices $t$ of $T$ such that the path of $T$ between $u$ and $t$ contains $v$. We show by induction on $j \ge 0$ that if every path of $T$ with first vertex $u$ and second vertex $v$ has at most $j$ vertices in $N(T)$ other than $u$, then the branch at $u$ through $v$ contains at most $2^j - 1$ members of $N(T)$. We may assume that the branch contains a member of $N(T)$; let $y$ be one for which the path $Q$ of $T$ between $u$ and $y$ is shortest. Then $j \ge 1$. The vertices of $Q$ other than $u, y$ are not in $N(T)$, and so have degree $2$ in $T$; hence every other member of $N(T)$ in the branch is in the branch at $y$ through a neighbour of $y$ not in $Q$, and $y$ has at most two such neighbours, since $T$ has maximum degree $\le 3$. For each such neighbour $y'$, every path with first vertex $y$ and second vertex $y'$ has at most $j - 1$ vertices in $N(T)$ other than $y$, since adding $Q$ to it gives a path with first vertex $u$ and second vertex $v$, in which $y$ and these vertices are in $N(T)$. By induction, each of these branches contains at most $2^{j-1} - 1$ members of $N(T)$, and so the branch at $u$ through $v$ contains at most $1 + 2(2^{j-1} - 1) = 2^j - 1$. Now if $|V(T)| \ge 2$, let $r$ be a vertex of $T$ of degree $1$, with neighbour $v$. Then $r \notin N(T)$, every other vertex of $T$ is in the branch at $r$ through $v$, and by (2) every path with first vertex $r$ and second vertex $v$ has at most $n$ vertices in $N(T)$. Hence $|N(T)| \le 2^n - 1$.}\fixnote{fix:12.3-count-proof-added}{The paper gives no argument for this bound. The argument above is the one used in the formalization.} Let $Z = \bigcup (V(W_t) : t \in N(T))$; then $|Z| \le 2^n(w' + 1)$. Now every component of $T \setminus N(T)$ is a path, and so we may add edges to $T \setminus N(T)$ to obtain a path $P$. For each $t \in V(P)$, let $W'(t) = W(t) \cap (G \setminus Z)$.

\hspace*{12pt}(3) \textit{$(P, W')$ is a path-decomposition of $G \setminus Z$.}
